\section{The canonical product, convergence, zeros, and entirety}\label{sec:product}

Put
\begin{equation}\label{eq:Fdef}
 F(z):=G(1+z)
 =(2\pi)^{z/2}
 \exp\!\left(-\frac{z+(1+\gamma)z^2}{2}\right)
 \prod_{k=1}^{\infty}E_k(z),
\end{equation}
where
\begin{equation}\label{eq:Ek}
 E_k(z):=\left(1+\frac zk\right)^k
          \exp\!\left(-z+\frac{z^2}{2k}\right).
\end{equation}
This is the product displayed by DLMF, Wikipedia and MathWorld.  Its form is
already suggestive of a constructive proof: after the linear and quadratic
terms of \(k\Log(1+z/k)\) are cancelled, the tail is summable like \(k^{-2}\).

\subsection{The one estimate that drives the construction}

For \(|z|\le R\) and \(k\ge 2R\), use the principal Taylor branch of the
logarithm in \(|z/k|\le 1/2\) and set
\[
 q_k(z):=k\Log\left(1+\frac zk\right)-z+\frac{z^2}{2k}.
\]
The first two Taylor terms cancel exactly:
\begin{equation}\label{eq:qseries}
 q_k(z)=\sum_{m=3}^{\infty}
       \frac{(-1)^{m+1}z^m}{m k^{m-1}}.
\end{equation}

\begin{lemma}[Uniform canonical-factor tail]\label{lem:canonical-tail}
If \(|z|\le R\), \(N\ge 2R\), and \(N\ge1\), then
\begin{align}
 |q_k(z)|&\le \frac{2R^3}{3k^2}\qquad(k>N),\label{eq:qbound}\\
 \sum_{k>N}|q_k(z)|&\le \frac{2R^3}{3N}.
 \label{eq:qtail}
\end{align}
Consequently, for \(M>N\),
\begin{equation}\label{eq:product-cauchy}
 \left|\prod_{k=N+1}^{M}E_k(z)-1\right|
 \le \exp(B_N)B_N,
 \qquad B_N=\frac{2R^3}{3N}.
\end{equation}
\end{lemma}

\begin{proof}
From \cref{eq:qseries} and \(R/k\le1/2\),
\[
 |q_k(z)|
 \le \sum_{m=3}^{\infty}\frac{R^m}{m k^{m-1}}
 \le \frac{R^3}{3k^2}\sum_{j=0}^{\infty}\left(\frac Rk\right)^j
 \le \frac{2R^3}{3k^2}.
\]
The elementary integral comparison
\[
 \sum_{k=N+1}^{\infty}\frac1{k^2}
 \le \int_N^{\infty}\frac{dt}{t^2}=\frac1N
\]
is proved arithmetically by the telescoping inequality
\(
1/k^2\le1/(k-1)-1/k
\)
for \(k\ge2\), after clearing denominators.  This proves
\cref{eq:qtail}.  Finally
\(
\prod E_k=\exp(\sum q_k)
\)
and
\(
|\e^w-1|\le \e^{|w|}|w|
\), whose finite Taylor proof has an explicit remainder.  Applying it to the
finite sum and then passing through the Cauchy modulus gives
\cref{eq:product-cauchy}.
\end{proof}

Every inequality in this proof has a division-free version.  A convenient
primitive-recursive cutoff is
\begin{equation}\label{eq:product-cutoff}
 N(R,p):=\max\!\left(2R,\left\lceil\frac{2^{p+5}R^3}{3}\right\rceil\right).
\end{equation}
Then \(B_N\le2^{-p-5}\).  On this range \(B_N\le1\), so
\(\e^{B_N}B_N\le3B_N\le2^{-p-3}\), leaving guard bits for the
finite factors and the prefactor.  The ceiling is implemented by bounded
integer division, and the verifier checks the cross-multiplied inequality
rather than trusting the formula for \(N\).

\subsection{Derivative tails}

Uniform convergence of values alone gives a continuous function.  To obtain a
constructive holomorphic contract, differentiate the logarithmic tail.  Direct
calculation gives
\begin{align}
 q_k'(z)&=\frac1{1+z/k}-1+\frac zk
        =\frac{(z/k)^2}{1+z/k},\label{eq:qprime}\\
 q_k''(z)&=-\frac{1}{k(1+z/k)^2}+\frac1k
        =\frac{2z/k+z^2/k^2}{k(1+z/k)^2},\label{eq:qsecond}\\
 q_k^{(r)}(z)&=\frac{(-1)^{r-1}(r-1)!}
 {k^{r-1}(1+z/k)^r}\qquad(r\ge3).
 \label{eq:qhigher}
\end{align}
Thus, when \(|z|\le R\) and \(k\ge2R\),
\begin{align}
 |q_k'(z)|&\le \frac{2R^2}{k^2},\label{eq:qprime-bound}\\
 |q_k''(z)|&\le \frac{12R}{k^2},\label{eq:qsecond-bound}\\
 |q_k^{(r)}(z)|&\le
 2^r(r-1)!\,k^{1-r}\qquad(r\ge3).
 \label{eq:qhigher-bound}
\end{align}
For each fixed derivative order the tails have explicit moduli.  At orders
\(r\ge3\), the geometric/integral estimate is even faster.  These bounds give
a finite Taylor-ball certificate on every rational disk: evaluate finitely many
factors and sum finitely many derivative tails.  No appeal to a nonconstructive
``uniform limit theorem'' is necessary; that theorem is re-proved here in the
only quantitative form required.

\begin{theorem}[Constructive existence and entirety]\label{thm:entire}
The partial products in \cref{eq:Fdef} form, on every rational compact disk, a
fast uniformly Cauchy family together with every fixed derivative.  Their limit
\(F\) has a primitive-recursive evaluation modulus and compatible Taylor-ball
certificates.  Hence \(G(z)=F(z-1)\) is entire in the constructive sense of
\cref{sec:ipra}.
\end{theorem}

\begin{proof}
Fix an integer radius \(R\).  The first \(2R\) factors and the prefactor are
finite combinations of polynomials and exponentials, for which the Gamma packet
already specifies rational-ball evaluators.  Apply
\cref{lem:canonical-tail,eq:qprime-bound,eq:qsecond-bound,eq:qhigher-bound} to
the remaining factors.  Given precision \(p\) and derivative order \(r\), a
primitive-recursive maximum of the corresponding cutoffs makes every tail at
most \(2^{-p-c_r}\).  The finite product rule and exponential error lemmas
combine them into a Taylor ball.  Compatibility of overlapping balls follows
from the same value-tail estimate.  All loops are bounded by the cutoff, so the
certificate producer is primitive recursive; all local checks are rational
inequalities, so its soundness proof uses quantifier-free induction.
\end{proof}

\subsection{Normalization and real conjugation}

At \(z=0\), every factor in \cref{eq:Fdef} is one and the prefactor is one.
Therefore
\begin{equation}\label{eq:G1}
 G(1)=F(0)=1.
\end{equation}
This is exact, not merely a limiting statement.

Every coefficient and constant in the product is real.  Finite partial products
therefore satisfy
\(
F(\overline z)=\overline{F(z)}
\)
exactly.  The common tail modulus permits passage to the named limits:
\begin{equation}\label{eq:conj}
 G(\overline z)=\overline{G(z)}.
\end{equation}
In iPRA the proof at precision \(p\) simply evaluates both sides using the same
cutoff; the finite centers are conjugate and the two radius bounds coincide.

\subsection{Zeros and their multiplicities}

\begin{theorem}[Zero divisor]\label{thm:zeros}
The zeros of \(G\) are exactly
\[
 0,-1,-2,\ldots,
\]
and the zero at \(1-m\) has multiplicity \(m\) for every \(m\ge1\).
\end{theorem}

\begin{proof}
For \(F(z)=G(1+z)\), the finite factor
\((1+z/m)^m\) vanishes to order \(m\) at \(z=-m\).  Its exponential multiplier
never vanishes.  Every other finite factor is nonzero there, as is the
prefactor.  Choose \(N>2m+2\).  On a sufficiently small rational disk around
\(-m\), all factors with \(k>N\) satisfy \(|z/k|<1/2\); their logarithmic tail
converges to a finite value, so their product is quantitatively apart from zero.
The remaining finitely many nonvanishing factors have a rational lower bound on
a smaller disk.  Therefore \(F(z)=(z+m)^m U_m(z)\), where \(U_m\) is analytic
and apart from zero on that disk.

Conversely, at a point that is not a negative integer, choose a rational disk
avoiding the finitely many possible zeros among the head factors.  Their product
has a positive rational lower bound, and the tail exponential is nonzero.  Thus
there is no other zero.  Translating by one gives the stated divisor of \(G\).
\end{proof}

The factorization proof is fully finitary: for a fixed integer \(m\), polynomial
division extracts \((z+m)^m\) from the \(m\)-th factor, while all lower bounds
are supplied by bounded searches over rational boxes.  This proves the integer
zero statements on both target pages and strengthens them by multiplicity.

\subsection{Why the product is the right iPRA definition}

One could instead define \(G\) from its functional equation plus a uniqueness
principle.  That route postpones rather than removes the analytic work: the
uniqueness condition involves third derivatives or an asymptotic modulus.  The
canonical product is preferable as the foundational iPRA definition because
its convergence is controlled by the single rational inequality
\cref{eq:qtail}.  Functional equations, reflection and special values then
become theorems with independently checkable finite certificates.
